\begin{abstract}
Source apportionment from sparse urban air-quality sensors is an inverse problem whose resolution is limited by sensor placement, wind-driven transport, background variation, and noise.  Known or proxy emission inventories make attribution meaningful by restricting the source field to a finite set of candidate groups, but do not guarantee that those groups are distinguishable from the available observations.  We formulate inventory-based apportionment as a wind-conditioned lagged inverse problem in which each source and temporal-basis coefficient produces a sensor-time fingerprint.  After projecting out a low-dimensional background space, separability is governed by the projected response matrix \(\widetilde H_\Phi\): exact identifiability requires its full column rank, and noise-robust attribution its singular values, coefficient visibility, background absorption, and pairwise coherence, so predictive fit alone is not evidence of identifiability.  Our identifiability-aware apportionment (IASA) framework fits nonnegative source--basis coefficients, reconstructs activity trajectories, reports uncertainty, and recommends conservative groupings for indistinguishable sources.  Instantiated on a New Delhi platform (government PM\(_{2.5}\) and wind records, regulatory sensor locations, and four proxy source groups) with controlled and observed-data studies, it reports the attribution resolution defensible under the declared inventories, transport, background, lag, and noise rather than the most detailed possible attribution vector.
\end{abstract}
